\chapter{Resultados e Discussões}

O desenvolvimento do \textit{software} demonstrou que é possível integrar conceitos fundamentais da engenharia química com tecnologias modernas de programação para produzir uma ferramenta acessível, modular e capaz de automatizar cálculos recorrentes da área.

Do ponto de vista funcional, o \textit{backend} em \textit{FastAPI} apresentou desempenho satisfatório, respondendo rapidamente aos cálculos mesmo em operações envolvendo chamadas reiteradas ao módulo de propriedades termodinâmicas (\textit{CoolProp}) ou a rotinas numéricas como métodos iterativos de Brent e integrações do \textit{SciPy}.

A validação de entradas pelo \textit{Pydantic} mostrou-se fundamental para evitar inconsistências numéricas e garantir segurança nos cálculos, reduzindo significativamente erros comuns como unidades incorretas, tipos inválidos ou parâmetros ausentes.

No \textit{frontend}, o uso de \textit{React}, \textit{TypeScript} e \textit{Vite}, associado a \textit{Tailwind CSS}, primitivas acessíveis de \textit{@base-ui/react}, componentes compostos no padrão \textit{shadcn}, \textit{React Router} e elementos visuais próprios, resultou em uma interface modular, responsiva e adequada ao propósito educacional do sistema. Essa arquitetura favoreceu a separação entre apresentação, estado e integração com a API, permitindo encapsular cada categoria de cálculo em componentes e serviços tipados. Bibliotecas como \textit{KaTeX}, \textit{Recharts}, \textit{Lucide React} e \textit{Sonner} complementam a interface com renderização matemática, gráficos reativos, ícones contextuais e notificações discretas. Além disso, a criação de uma camada compartilhada para grades, eixos e rótulos numéricos padronizou parte relevante das visualizações em \textit{SVG}, o que melhora a consistência gráfica e simplifica a evolução dos componentes. Na interface consolidada, a navegação evidencia o papel pedagógico de cada área, com uma home orientada por trilhas, um núcleo de simulações mais explícito e indicadores visuais de processo que reforçam a interpretação do regime de escoamento.

Para além da automação de cálculos, a incorporação de recursos didáticos consolidou o DCOU como um ambiente de ensino, e não apenas como uma calculadora. As explicações teóricas embutidas em cada módulo, o glossário técnico e o conjunto ampliado de visualizações interativas aproximam o resultado numérico do conceito que o fundamenta, ao apresentar o comportamento dos fenômenos físicos de forma contextualizada. Os exercícios integrados, por sua vez, mostraram-se eficazes em articular múltiplos módulos em um único fluxo de trabalho, reproduzindo a sequência de decisões de um problema real de engenharia e evidenciando como os cálculos isolados se conectam em uma análise completa. No estado final do projeto, esses recursos passaram a coexistir com um \textit{shell} de navegação persistente, cartões de entrada rápida, indicadores compactos e blocos visuais voltados explicitamente a docentes e discentes, reforçando a utilidade pedagógica da plataforma.

A comparação dos resultados obtidos com valores de referência da literatura e com \textit{softwares} consolidados evidenciou boa concordância para todos os módulos implementados. Cálculos de \textit{Reynolds}, perdas de carga, seleção de diâmetros e estimativas de \textit{NPSH} apresentaram diferenças mínimas em relação às tabelas clássicas. Para reatores CSTR e PFR, os valores de conversão e volume mostraram-se consistentes com soluções de Fogler e Levenspiel, reforçando a confiabilidade das integrações numéricas.

Embora o \textit{software} esteja plenamente funcional para os módulos desenvolvidos, observou-se que o sistema ainda pode ser expandido significativamente. A ausência de módulos como trocadores de calor, destiladores, absorção, secagem e operações com múltiplas reações limita sua aplicabilidade em projetos de maior complexidade. Do ponto de vista de usabilidade e de maturidade funcional, futuras melhorias incluem exportação de relatórios, ampliação de gráficos especializados e integração com bancos de dados externos de propriedades físicas, de modo a ampliar o alcance técnico da ferramenta.

Ainda assim, o objetivo de criar uma ferramenta precisa, acessível, gratuita e colaborativa foi alcançado. O fato de o projeto ser \textit{open-source} torna natural que seu ciclo de evolução dependa da contribuição da comunidade acadêmica e profissional, que poderá incluir novos módulos, otimizar algoritmos e expandir a base de dados.
